\subsection{Certification of the fixed-depth lifting theorem}

\begin{result}{C.2.12}{Proved}{SECTIONS A.65, with its arithmetic and localization inputs proved}
\label{thm:ghost-fixed-depth-survival}
For every integer \(R\ge1\) and positive integer \(N\),
there is an ordinary infinite-support completed pair
\((A,B)\) and a continuous section
\[
 s:G_{\mathbb Q}\longrightarrow
 E_R=\pi_1(U,\overrightarrow{01})/\Gamma_{R+1}\Pi
\]
whose abelianization is \((A,B)\). Here all lower
central subgroups are closed. The section agrees with
the rational-point section of three in the actual
finite arithmetic quotient defined by truncating the
integral Magnus expansion at degree \(R\) and
reducing its coefficients modulo \(N\), as specified
in the proof. Its abelian pair is neither a rational
point image nor a cuspidal image. The pair and the
section in this statement are allowed to depend on
\(R,N\).
\end{result}
\begin{proof}
We prove the assertion for \(R\ge3\); projection
gives the two smaller classes. Fix a representative
\(s_3\) of the rational-point section of three and
use its actual conjugation action \(\rho_3\) to
write \(E_r=N_r\rtimes_{\rho_3}G_{\mathbb Q}\),
where \(N_r=\Pi/\Gamma_{r+1}\Pi\). Thus a relative
geometric cocycle \(z\) gives the section
\(g\mapsto z(g)s_3(g)\). For each coefficient prime
the Hall model is Lemma~\ref{lem:ghost-hall-model}.
Every finite nilpotent group is the direct product
of its Sylow subgroups, so the universal finite
nilpotent quotients identify \(N_r\) with the
product of those pro-\(\ell\) Hall models. In
particular its central layer is
\[
 L_r=\Gamma_r\Pi/\Gamma_{r+1}\Pi
       =\widehat{\mathbb Z}(r)^{b_r},
 \tag{FD.1}
\]
where \(b_r\) is the number of weight-\(r\) basic
commutators. The Tate action follows because the
action on both generators modulo commutators is
cyclotomic and an inner automorphism is trivial
on this central layer.

\emph{The exact finite congruence.}
Let \(I=(X,Y)\) in the noncommutative Magnus
algebra. Define the characteristic open subgroup
\[
 K_r(m)=\ker\left(N_r\longrightarrow
       (1+I)\bigl((\mathbb Z/m)\langle X,Y\rangle/I^{r+1}\bigr)
                       \right).
 \tag{FD.2}
\]
The map sends the two geometric generators to
\(1+X,1+Y\). Substitution by the integral Magnus
series of any two generator images preserves the
ideal of coefficients divisible by \(m\); hence
\(K_r(m)\) is stable under every continuous
automorphism. If \(\bar\rho_3\) is the induced
finite action, the actual arithmetic quotient is
\[
 E_r\longrightarrow
 (N_r/K_r(m))\rtimes\operatorname{im}\bar\rho_3,
 \qquad (z,g)\longmapsto(z\bmod K_r(m),\bar\rho_3(g)).
 \tag{FD.3}
\]
Agreement of sections means \(z(g)\in K_r(m)\)
for every \(g\), which is equality of their maps
under (FD.3).

Let \(h_R\) be the number of Hall commutators of
weights at most \(R\), and put
\[
 D=(R!)^{h_R+1},\qquad T=2(R!)^R.
 \tag{FD.4}
\]
The continuous set lift \(\sigma_r:N_{r-1}\to N_r\)
sets all new Hall coordinates to zero. It satisfies
\[
 \sigma_r(K_{r-1}(Dm))\subseteq K_r(m).
 \tag{FD.5}
\]
Indeed the lowest Magnus polynomials of the Lyndon
Hall words have distinct leading words with
coefficient one. Their triangular extraction
matrices are integral. Extract a Hall exponent
and remove its power, successively. If an
exponent \(a\) belongs to an integral ideal \(J\),
every nonconstant coefficient of its Hall power
through degree \(R\) belongs to \(J/(R!)\),
using
\(\binom ak=(a/k)\binom{a-1}{k-1}\), \(k\le R\).
There are at most \(h_R\) removals. The additional
factor \(R!\) in \(D\) lets each extracted power
be extended through degree \(r\) with all its
coefficients divisible by \(m\). Multiplying
these extended powers proves (FD.5). This
argument works in every \(\mathbb Z_\ell\),
including coefficient primes dividing \(R!\).

\emph{A uniform annihilator for the retained global kernel.}
For \(2\le r\le R\), set
\[
 B_r=2\lvert K_{2r-2}(\mathbb Z)\rvert,
 \qquad D_r^{\rm old}=(2^{r-1}-1)(3^{r-1}-1).
 \tag{FD.6}
\]
These are positive finite integers. Finite
generation is Quillen's theorem and the even
rank is zero by Borel's theorem
\cite[IV.6.9, IV.1.18]{GhostWeibelKbook}.
We have, for every coefficient prime,
\[
 B_rH^2_{\rm cts}
 (G_{\mathbb Q,\{\ell,\infty\}},\mathbb Z_\ell(r))=0.
 \tag{FD.7}
\]
For odd \(\ell\), the comparison in
\cite[VI.8.2]{GhostWeibelKbook} identifies this
group with the \(\ell\)-primary part of
\(K_{2r-2}(\mathbb Z)\). Localization at \(\ell\)
does not change that primary part: the adjacent
positive finite-field \(K\)-groups have orders
prime to \(\ell\). At two apply VI.9.11 to
\(\mathbb Z[1/2]\). The real rank is one, so
the integer \(\rho\) there is zero. For
\(r\not\equiv0\pmod4\) the group is the
corresponding two-primary \(K\)-group. For
\(r\equiv0\pmod4\), its kernel under real
restriction is that group and its remaining
image has exponent two. The extra factor two
in (FD.6) gives (FD.7) in this case also.
The comparison of continuous Galois cohomology
with this integral etale cohomology is the
coefficient-compatible comparison used in
Theorem~\ref{thm:ghost-strict-localization}.

\emph{Control of actual inertia and of its lifted relator.}
Fix \(q\ne\ell\). The tangential action is
unramified in every nilpotent class. In class
two this follows from the actual Ihara action
in Corollary~\ref{cor:ghost-class2-lift} and
the unramified cyclotomic character. Inductively,
the kernel of the action map between consecutive
classes is a submodule of
\(\operatorname{Hom}(N_1,L_{i+1})\), with
Frobenius weight \(i\). Wild and non-\(\ell\)
tame inertia have zero image in this pro-\(\ell\)
module. The remaining tame generator has an
image killed by \(q^i-q\), comparing Frobenius
conjugation with \(FTF^{-1}=T^q\). For
\(i\ge2\) this is a nonzero integer acting on
a torsion-free module, so the image is zero.
This proves the induction. At \(q>3\),
\(s_3\) is unramified as well: its two Kummer
coordinates are units, and each successive
central inertia coordinate is killed by
\(q^i-q\), starting at weight two. Thus the
relative action is unramified at all fresh
support primes.

At such a prime write an ordinary valuation
pattern as \((E,0),(0,E)\), or \((-E,-E)\).
Every local integral class-two section has
tame inertia \(V_2=y^n x^m C^t\) with
\[
 m,n\in E\mathbb Z_\ell,
                 \qquad t\in(E/2)\mathbb Z_\ell.
 \tag{FD.8}
\]
For the positive patterns the local primitive
calculation in Proposition~\ref{prop:ghost-ordinary-cup}
gives \(t\) equal, up to sign, to the nonzero
valuation times an integral primitive. In the
negative pattern it is the sum of a binomial
expression in that valuation and its product
with an integral primitive, hence belongs to
\((E/2)\mathbb Z_\ell\). These statements
apply to every local lift: the difference
of two such central cochains is a local
\(\mathbb Z_\ell(2)\)-cocycle, whose inertia
is killed by \(q-1\), hence is zero.
Since \(s_3(T)=1\), the same inertia
coordinates occur for the relative cocycle.
Its canonical Hall extension through degree
\(R\) is therefore one modulo the Magnus
coefficient ideal \(E/(2R!)\).

Suppose a section through class \(i\) is
given. Let \(\bar V\) be the canonical lift
of its inertia from class \(i-1\), and write
its actual last coordinates as \(Z_i\in L_i\).
For its actual Frobenius \(U\) and the
actual unramified action \(\varphi_q\),
the tame relation in class \(i\) is
\[
 U\varphi_q(\bar V)U^{-1}\bar V^{-q}
                  =Z_i^{-q(q^{i-1}-1)}.
 \tag{FD.9}
\]
Integral substitution, conjugation and powers
preserve any integral Magnus coefficient ideal.
The left side is central, so its weight-\(i\)
Hall coordinates are extracted by the integral
triangular matrix, with no denominator.
Division by \(q^{i-1}-1\) bounds \(Z_i\);
extending its Hall powers through degree \(R\)
loses at most a further factor \(R!\).
Starting with (FD.8) proves by induction that
the canonical lift of inertia through class
\(i\) is one modulo
\[
 \frac{E}{2(R!)^{i-1}\prod_{k=2}^{i-1}(q^k-1)}
                                      \mathbb Z_\ell
                    \quad(2\le i\le R).
 \tag{FD.10}
\]
Every ideal used below is integral because
of the eventual choice of \(E\).

For a canonical set lift from class \(r-1\)
to \(r\), let \(c\) be its normalized central
defect. Its exact strict residue is
\(z_q(c)=q^{-1}\) times the lifted tame
relator, by Proposition~\ref{prop:ghost-tame-cochain}.
Consequently (FD.10) gives
\[
 \frac{z_q(c)}{q^{r-1}-1}
 \in
 \frac{E}{T\prod_{k=2}^{r-1}(q^k-1)}
                         \mathbb Z_\ell^{b_r}.
 \tag{FD.11}
\]
For \(r\ge3\) use (FD.10) in class \(r-1\)
and divide its relator once more. For \(r=2\),
an ordinary local class-two lift exists by
the explicit local primitives already used.
Its last inertia coordinate (FD.8) corrects
the canonical relator, giving the stronger
bound \(z_q(c)/(q-1)\in(E/2)\mathbb Z_\ell\).
The product in (FD.11) is empty in this case.
When both valuations are zero, successive
torsion-free tame relations make every
inertia coordinate zero; the canonical
defect then has residue literally zero.

All defects here belong to the actual tame
cochain domain of Proposition~\ref{prop:ghost-tame-cochain}.
Indeed the tangential action is unramified,
and any pro-\(\ell\) section kills wild and
non-\(\ell\) tame inertia. Its relative
cocycle and its Hall set lift therefore
factor through the indicated tame quotient,
including at two and three when \(q\ne\ell\).
This justifies using (FD.11) as a representative
of the global strict localization map.

\emph{Induction with the finite congruence retained.}
Put \(M_R=N\) and, backwards, define
\[
 M_{r-1}=DB_rM_rD_r^{\rm old}
                         \quad(2\le r\le R).
 \tag{FD.12}
\]
Choose \(C\) divisible by \(24,M_1\) and
all \(TB_rM_r\). Apply
Theorem~\ref{thm:ghost-general-exponent-arithmetic}
with \(d=R-1\). It gives one ordinary
pair whose fresh exponents are
\[
 E_q=C e_q\prod_{k=2}^{R-1}(q^k-1),
              \qquad e_q\longrightarrow0
                         \text{ in }\widehat{\mathbb Z}.
 \tag{FD.13}
\]
Its quotients \(A/3,B/(-2)\) are global
\(M_1\)-th powers, so choose literal relative
abelian cocycles pointwise divisible by
\(M_1\). This starts the section induction.

Given a section through class \(r-1\)
whose relative cocycle lies pointwise in
\(K_{r-1}(M_{r-1})\), use \(\sigma_r\).
By (FD.5) its set lift lies pointwise in
\(K_r(L)\), where
\(L=B_rM_rD_r^{\rm old}\). Its multiplication
defect has Magnus coefficients divisible
by \(L\); central Hall extraction is
integral, so its normalized central
two-cocycle \(c\) has values in \(LL_r\).
Thus \(c/(B_rM_r)\) is an actual integral
continuous two-cocycle.

At two and three away from \(\ell\),
its strict residue is divisible by
\(D_r^{\rm old}\), and hence by the
needed \(q^{r-1}-1\). This follows
directly from the integral cochain
formula for \(z_q\); division by \(q\)
is by an \(\ell\)-adic unit. At each
fresh prime, (FD.11), (FD.13), and the
choice of \(C\) put the divided residue
in \(e_q\mathbb Z_\ell^{b_r}\). At every
other prime it is zero. These quotients
are integral and tend to zero, so belong
to the literal module \(R_\ell^{b_r}\)
of Theorem~\ref{thm:ghost-strict-localization}.
That theorem places
\[
                 \beta=[c/(B_rM_r)]
 \tag{FD.14}
\]
in the image of the unramified global
group at every coefficient prime.
Equation (FD.7) gives \(B_r\beta=0\).
Therefore \([c]=0\), and there is a
continuous integral one-cochain \(h\)
with \(\delta h=-c\), producing an
actual class-\(r\) section.

It remains to retain the specified
congruence. Reducing \(h\) modulo
\(B_rM_r\) gives a one-cocycle, and
the coefficient connecting maps satisfy
\[
 \partial_{B_rM_r}[\bar h]=-\beta,
 \qquad
 \partial_{M_r}[\bar h\bmod M_r]
                    =-B_r\beta=0.
 \tag{FD.15}
\]
The second equality is the literal
identity \(\delta h/M_r
       =B_r(\delta h/(B_rM_r))\).
The long exact sequence for
\(0\to L_r\xrightarrow{M_r}L_r\to L_r/M_r\to0\)
therefore lifts the reduced cohomology
class to an integral one-cocycle \(z\).
If its representative differs from
\(h\) modulo \(M_r\) by a coboundary,
lift the finite zero-cochain and adjust
\(z\) by its integral coboundary.
Then \(z\equiv h\pmod{M_r}\) pointwise.
Replacing \(h\) by \(h-z\) leaves
\(\delta h=-c\) and makes its new
coordinates divisible by \(M_r\).
Together with the set lift this gives
agreement in \(K_r(M_r)\), completing
the induction.

All coefficient sequences used here
are sequences of continuous cochains:
a finite-valued cochain lifts by lifting
each of its finitely many values, so
the asserted long exact sequences have
their usual cochain derivation. Cochains
into a product of compact modules are
exactly products of component cochains,
and their differentials are componentwise.
The corrections for all \(\ell\)
therefore combine into one continuous
section of the full integral quotient.
At the last step \(M_R=N\), giving
the exact finite agreement (FD.3).

Finally, the arithmetic theorem's
valuation display (VX.8)
gives infinitely many nonzero valuation
coordinates. A rational point has finite
support. The actual abelian cusp maps
have second coordinate trivial at zero,
first coordinate trivial at one, and
equal valuation coordinates at infinity.
The fixed valuation pairs at two and
three exclude each of these maps.
This proves all assertions of A.65 with
the omitted arithmetic input reconstructed.
\end{proof}
